\subsection{Terminals through twenty-four}

\begin{proposition}
\label{prop:small-terminals}
Theorem~\ref{thm:main} holds for \(2\le N\le24\).
\end{proposition}
\begin{proof}
Lemma~\ref{lem:constant-support} covers \(a=0\). At \(N=2\), \(R_1\)
has a real root and \(R_2\) is positive on \(\R\) by its gamma integral.
Take \(a\ge1\) and \(3\le N\le24\). The weak annulus in
Lemma~\ref{lem:gamma-annulus} gives \(a<d(N-2)\) if the two polynomials
share a root.

Write \(s=d^2/(a+d)\) and use the endpoint normalization
of Lemma~\ref{lem:center-ratio}. If \(s\ge20\), the mixed
coefficients at either order have sum less than \(1/2\).
Indeed their upper bound in \eqref{eq:normalized-coefficients}
is symmetric, and for \(k\le m/2\) the exponential term is
at most \(e^{-5k}\). Thus the sum is bounded by
\[
 2\bigl((1+e^{-5})^{24}-1\bigr)
 <2\bigl(e^{6/35}-1\bigr)<\frac{12}{29}<\frac12.
\]
Here \(e^{-5}<1/140\), and \(e^x<1/(1-x)\) for \(0<x<1\).
In the normalized variable the roots have moduli between
\(2^{-1/m}\) and \(2^{1/m}\), by comparison of \(1+w^m\)
with the mixed terms and by reversal. Since \(d>s\ge20\),
\[
 \frac d2\log\!\left(\frac{N}{N-1}\right)
 \frac{N(N-1)}{2N-1}
 \ge \frac{d(N-1)}{2(2N-1)}\ge 4>\log 2 .
\]
The center ratio in Lemma~\ref{lem:center-ratio} therefore
separates the annuli. A common root requires \(s<20\).

We next show that \(d<43\). If \(d\ge43\), set
\(r=\lfloor d/43\rfloor\). Then \(43\le d/r<86\), and the interval
\((aN/r,(aN+2d)/r]\) has length at least \(86\). Its upper endpoint
is below \(86[N(N-2)+2]\le45580\).
The bound \(s<20\) gives \(a>d^2/20-d\ge d^2/43\ge rd\), so the
left endpoint exceeds \(dN\). The first record prime gap larger
than \(86\) starts at \(360653\) \cite[Table~8]{oliveira-gaps}.
It follows that the interval contains a prime \(p>dN\), and that
\(rp\) belongs to the first \(2d\) factorial terms at order \(N\).

Since \(a<d(N-2)\), we have \(p^2>(a+d)N\).
Each factorial interval \((am,(a+d)m]\) contains at most one
multiple \(u_m\) of \(p\), and all binomial coefficients are
\(p\)-adic units. A positive valuation of a reciprocal root
has the form \(1/w_m\), where
\(w_m=1+\lfloor((a+d)m-u_m)/d\rfloor\).
Equal positive valuations at both orders would imply
\(u_N-u_{N-1}\in[a+1,a+2d-1]\subset(0,p)\), which is impossible.
Thus all roots of a common factor would be \(p\)-adic units.
But the first \(2d\) terms at order \(N\) contain \(rp\), so
\(R_N\bmod p\) has at most one root away from zero, counted with
multiplicity. A monic common divisor reduces with its degree
unchanged, and would have nonzero constant term and even degree
at least two. This contradiction proves \(d<43\).

It remains to test the integer triples
\[
 3\le N\le24,\qquad 1\le d\le42,\qquad
 1\le a<d(N-2),\qquad d^2<20(a+d).
\]
The supplementary verifier exhausts this set and finds
\(217\,952\) triples. At each triple it forms the two monic
reciprocals modulo \(1\,000\,000\,007\), a prime certified by
trial division, and computes their gcd by Euclidean division.
Every gcd is one. A monic integral common divisor would remain
a nonconstant common divisor on reduction. This proves the proposition.
\end{proof}
